\section{Qualitative Analysis}

The quantitative pattern is easy to see in individual scenes. In scene \texttt{ps\_000005}, the target and a distractor are both small green spheres. Mirror self-play selects ``It's a small green sphere.'' Its local partner resolves this correctly, but two of the three alternate-model held-out listeners choose the distractor. Population-play selects ``The target is the small green sphere at row 4, column 1,'' which all held-out listeners map to the target. The mirror message is syntactically valid but pragmatically under-informative.

In scene \texttt{ps\_000011}, mirror self-play selects ``Find the large yellow cube left of the purple cylinder.'' Under perspective shift, this relation is frame-sensitive: one alternate-model held-out listener chooses the same-shape, same-color distractor. Population-play selects ``Go to row 4, column 1: the large yellow cube,'' which all held-out listeners map to the target. The failure is not merely brevity; it is an unshared orientation convention.

These examples show two ways same-partner success can mislead: the training partner may resolve ambiguity with a target-favoring tie-breaker, or it may share the speaker's spatial frame. Both produce perfect same-play success without robust communicative success.

The rubric-coded mirror failures show a concentrated failure mode. Across 152 listener-level failures, 147 are underspecified distractor cases and 5 are perspective-frame errors; no coded failures are wrong-attribute, private-landmark, or listener-misparse cases. A rule-based ambiguity verifier using visible attributes, row/column mentions, and left/right frame cues recovers all 152 coded mirror-failure labels, giving a non-LLM taxonomy check. A generated qualitative appendix adds two partial-observability examples, including a no-coordinate consensus+info repair that raises held-out success from 0.33 to 1.00 in the selected scene.

\section{Discussion}

The experiments support a modest but important conclusion: communicative success is not a sufficient evaluation signal unless the partner distribution is specified. Mirror self-play often has a working message available, but selects one that is under-specified for another partner. Population-play improves robustness by satisfying multiple listener assumptions, while the no-coordinate ablation shows that population diversity must be paired with a selector that values discriminative information.

This has implications for interactive language-agent research. Pipelines that generate alternatives, simulate a partner, and keep the response that appears to succeed can reward private conventions, hidden tie-breakers, or brittle spatial assumptions if evaluation uses the same partner type. Cross-play offers a low-cost diagnostic: hold out partner prompts, model families, or humans, and report the same-play/cross-play gap.

The no-coordinate ablation is also instructive. Exact coordinates are listener-invariant in our text grid, but may be unavailable or unnatural in embodied settings. The GPT-5.5 K=8 audit shows that robust non-coordinate expressions can be generated reliably here, while selector heuristics are not monotonic: adding more candidates improves population-play but hurts consensus+info. Robust communication requires both listener-visible evidence and selectors that choose genuinely discriminative descriptions.

\section{Limitations}

The current experiments are intentionally small and diagnostic. The selector partners are local deterministic listeners rather than trained LLM agents, so the result should be framed as a cross-play diagnostic rather than as evidence about large-scale agent training. The held-out listeners are API language models, not humans. We audit three held-out model families including \texttt{gpt-5.5}, but they are still language-model listeners and prompt variants rather than independent human partners. The strongest ceiling results use exact coordinate fallbacks, and the K=8 no-coordinate audit is limited to 50 perspective-stress scenes; it supports the generation-versus-selection diagnosis but not a broad claim about natural language grounding. The grid scenes are text-rendered, so the benchmark tests situated reference rather than embodied perception. Finally, the strongest evidence is concentrated in two diagnostic families, perspective shift and partial observability; broader claims require larger balanced runs across all scenario types and, ideally, human listener validation.

\section{Conclusion}

Same-partner communicative success is not sufficient evidence of robust pragmatic communication. Mirror self-play can achieve perfect same-play success while failing under held-out listeners, especially when messages are perspective-sensitive or under-specified. Population-play mitigates this when listener-invariant fallbacks are available, and the no-coordinate ablation shows that cross-play also diagnoses weaknesses in the selector itself. These results support cross-play as a simple evaluation axis for language agents that learn, filter, or select messages through situated interaction.
